\documentclass[11pt]{article}
\usepackage{mathblatt}
\usepackage{multicol}


\begin{document}
\pfheftstil
\pfheftkopf{Grundwert $G$ \textperiodcentered{} Lösungen}{}
{\small\noindent\textbf{Typische Fehler}\par\smallskip\noindent \textbullet\ p \% vom Teil gerechnet: 25 \% von 200 € statt 200 € · 4; 20 \% von 6 €.\par\noindent \textbullet\ Teil und Ganzes vertauscht.\par\noindent \textbullet\ falsches Ganzes (Teil zum Rest statt zum Ganzen).\par\noindent \textbullet\ „um“ und „auf“ verwechselt (um 30 \% gesenkt $\leftrightarrow$ auf 70 \% gesenkt).\par}\medskip
\begin{pfloesung}{}
\lz{1.}{1 \% = $\tfrac{1}{100}$ = 0,01; 10 \% = $\tfrac{1}{10}$ = 0,1; 20 \% = $\tfrac{1}{5}$ = 0,2; 25 \% = $\tfrac{1}{4}$ = 0,25; 50 \% = $\tfrac{1}{2}$ = 0,5; 100 \% = 1 = 1}{}{}
\end{pfloesung}
\begin{pfloesung}{}
\lz{2.}{der alte Preis, $60\,€$}{}{}
\end{pfloesung}
\begin{pfloesung}{}
\lz{3.}{größer; 70 € (21 € : 0,3)}{}{}
\end{pfloesung}
\begin{pfloesung}{}
\lz{4.}{$5$ Abschnitte bis $100\,\%$}{}{}
\end{pfloesung}
\begin{pfloesung}{}
\lz{5.}{$240$ kg}{wie oft bis hundert Prozent: $100 : 10 = 10$; mal nehmen: $10 \cdot 24 = 240$ kg}{}
\end{pfloesung}
\begin{pfloesung}{}
\lz{6.}{$800\,€$}{$1\,\% = 8\,€$; $100\,\%$}{}
\end{pfloesung}
\begin{pfloesung}{}
\lz{7.}{$225$ kg}{$81 : 36 = 2{,}25$; $2{,}25 \cdot 100$}{}
\end{pfloesung}
\begin{pfloesung}{}
\lz{8.}{800 €}{G = W : p = 200 € : 0,25}{}
\end{pfloesung}
\begin{pfloesung}{}
\lz{9.}{30 €}{G = W : p = 6 € : 0,2}{}
\end{pfloesung}
\begin{pfloesung}{}
\lz{10.}{450 € (360 € sind 80 \%)}{}{}
\end{pfloesung}
\begin{pfloesung}{}
\lz{11.}{G = 4 : $\tfrac{2}{3}$ = 6 Kugeln ($\tfrac{1}{3}$ sind 2 Kugeln)}{}{}
\end{pfloesung}
\begin{pfloesung}{}
\lz{12.}{G = 8 cm : 0,93 $\approx$ 8,60 cm}{}{}
\end{pfloesung}
\begin{pfloesung}{}
\lz{13.}{12,5 \%}{}{}
\end{pfloesung}
\begin{pfloesung}{}
\lz{14.}{110 €}{}{}
\end{pfloesung}
\end{document}